\begin{lemma}[The rotation average of the two real $L^1$ norms]
Let $v \in L^2(\gamma_n;\mathbb{C})$.  For $\theta \in \mathbb{R}$ put
$c_\theta = e^{i\theta}$ and let
\[
  f_1(\theta) = \operatorname{Re}(c_\theta v), \qquad
  f_2(\theta) = \operatorname{Im}(c_\theta v)
\]
be the real and imaginary parts of the rotated function, that is, the
$\gamma_n$-a.e.\ classes obtained by composing $c_\theta v$ pointwise with the
bounded $\mathbb{R}$-linear maps $\operatorname{Re},\operatorname{Im} :
\mathbb{C} \to \mathbb{R}$; they are elements of $L^2(\gamma_n;\mathbb{R})$,
since composition with a bounded linear map preserves membership in a Lebesgue
space.  Set
\[
  S_\theta = \|f_1(\theta)\|_{L^1(\gamma_n)} + \|f_2(\theta)\|_{L^1(\gamma_n)} ,
\]
the two $L^1$ norms being taken through the inclusion of \noderef{ToL1}.  Then:
\begin{enumerate}
\item the function $\theta \mapsto S_\theta$ is integrable on $[0,2\pi]$; and
\item
\[
  \frac{1}{2\pi}\int_0^{2\pi} S_\theta \, d\theta
    = \frac{4}{\pi}\int_{\mathbb{R}^n} |v(x)| \, d\gamma_n(x) .
\]
\end{enumerate}

This is the reference paper's displayed rotation average
$\frac{1}{2\pi}\int_0^{2\pi}S_\theta\,d\theta = \frac{4}{\pi}\|f\|_1$, with
$S_\theta = \|\operatorname{Re}(e^{i\theta}f)\|_1 + \|\operatorname{Im}(e^{i\theta}f)\|_1$,
and it is the single place the constant $8/\pi$ of the complex estimate comes
from: the $4/\pi$ here is $2 \cdot \frac{2}{\pi}$, one $\frac{2}{\pi}$ for the
cosine and one for the sine.  The right-hand side is written as the integral
$\int|v|\,d\gamma_n$ rather than as a norm in a complex $L^1$ space; this is the
same real number as the paper's $\|f\|_{L^1(\gamma;\mathbb{C})}$, and writing it
this way avoids introducing a complex $L^1$ inclusion that the tablet does not
yet have.  Clause (1) is not decoration: it is what Tonelli's theorem supplies
alongside the identity, and it is what licenses integrating an inequality in
$\theta$ against the constant function, which is how the average is used.
\end{lemma}
\begin{proof}
Write $M = \int_{\mathbb{R}^n}|v(x)|\,d\gamma_n(x)$ and fix a representative of
$v$, so that $v : \mathbb{R}^n \to \mathbb{C}$ is measurable and square
$\gamma_n$-integrable.

\emph{Step 0: $M$ is finite.}
By \noderef{GaussIsProb} the measure $\gamma_n$ is a probability measure, so the
constant function $1$ lies in $L^2(\gamma_n)$; by the integral form of
Cauchy--Schwarz (shared context) applied to $|v|$ and $1$,
$M \leq \bigl(\int |v|^2 d\gamma_n\bigr)^{1/2}\cdot 1 = \|v\|_{L^2} < \infty$.

\emph{Step 1: the one-variable integral of $|\cos|$ over a period.}
We claim $\int_0^{2\pi}|\cos\theta|\,d\theta = 4$.  Split the interval at
$\pi/2$ and $3\pi/2$, and use throughout that $\sin' = \cos$, so that the
fundamental theorem of calculus gives $\int_a^b \cos = \sin b - \sin a$ for all
$a \leq b$.  On $[0,\pi/2]$ we have $\cos \geq 0$, so
$\int_0^{\pi/2}|\cos| = \int_0^{\pi/2}\cos = \sin(\pi/2) - \sin 0 = 1$.  On
$[\pi/2, 3\pi/2]$ we have $\cos \leq 0$, so
$\int_{\pi/2}^{3\pi/2}|\cos| = -\int_{\pi/2}^{3\pi/2}\cos
 = -(\sin(3\pi/2) - \sin(\pi/2)) = -(-1-1) = 2$.  On $[3\pi/2, 2\pi]$ we have
$\cos \geq 0$ again, so
$\int_{3\pi/2}^{2\pi}|\cos| = \sin(2\pi) - \sin(3\pi/2) = 0 - (-1) = 1$.  The
integrand is continuous, hence integrable on each piece, and adding the three
pieces gives $1 + 2 + 1 = 4$.

\emph{Step 2: the same integral after a phase shift.}
For every $p \in \mathbb{R}$,
\[
  \int_0^{2\pi}\bigl|\cos(\theta + p)\bigr|\,d\theta = 4,
  \qquad
  \int_0^{2\pi}\bigl|\sin(\theta + p)\bigr|\,d\theta = 4 .
\]
For the first: substituting $\theta \mapsto \theta - p$ turns the integral into
$\int_p^{p+2\pi}|\cos u|\,du$, and $|\cos|$ is $2\pi$-periodic (because
$\cos(u + 2\pi) = \cos u$), so the integral of $|\cos|$ over any interval of
length $2\pi$ is the same as over $[0,2\pi]$, which is $4$ by Step 1.  For the
second: $\sin s = \cos(s - \pi/2)$, so
$|\sin(\theta+p)| = |\cos(\theta + (p - \pi/2))|$ for every $\theta$, and the
first identity applies with the phase $p - \pi/2$.

\emph{Step 3: the pointwise rotation identity.}
Fix $x \in \mathbb{R}^n$ and put $\rho = |v(x)| \geq 0$.  If $\rho = 0$ then
$v(x) = 0$ and both sides of the identity below are $0$.  Otherwise write
$v(x) = \rho e^{i\varphi}$ with $\varphi = \arg v(x)$, so that
$c_\theta v(x) = \rho e^{i(\theta + \varphi)}$ and therefore
\[
  \operatorname{Re}\bigl(c_\theta v(x)\bigr) = \rho\cos(\theta + \varphi),
  \qquad
  \operatorname{Im}\bigl(c_\theta v(x)\bigr) = \rho\sin(\theta + \varphi) .
\]
Taking absolute values, using $\rho \geq 0$, and integrating in $\theta$ with
Step 2 at the phase $p = \varphi$,
\begin{equation*}
  \int_0^{2\pi}\Bigl(\bigl|\operatorname{Re}(c_\theta v(x))\bigr|
      + \bigl|\operatorname{Im}(c_\theta v(x))\bigr|\Bigr)d\theta
    = \rho \cdot 4 + \rho \cdot 4 = 8\,|v(x)| . \tag{1}
\end{equation*}
In the degenerate case $\rho = 0$ the two sides of $(1)$ are $0$ as well, so
$(1)$ holds for every $x$.

\emph{Step 4: joint measurability.}
Let
$H(\theta, x) = \bigl|\operatorname{Re}(c_\theta v(x))\bigr|
  + \bigl|\operatorname{Im}(c_\theta v(x))\bigr| \geq 0$
on $[0,2\pi] \times \mathbb{R}^n$.  The map $(\theta,x)\mapsto c_\theta = e^{i\theta}$
is continuous, hence measurable; the map $(\theta,x)\mapsto v(x)$ is measurable,
being the composition of the measurable $v$ with the projection; their product is
measurable; and $\operatorname{Re}$, $\operatorname{Im}$ and $|\cdot|$ are
continuous.  So $H$ is measurable for the product $\sigma$-algebra, and it is
nonnegative.

\emph{Step 5: Tonelli.}
Lebesgue measure restricted to $[0,2\pi]$ is finite and $\gamma_n$ is a
probability measure by \noderef{GaussIsProb}, so both are $\sigma$-finite and
Tonelli's theorem (shared context) applies to $H$: the two iterated integrals
agree in $[0,+\infty]$.  Computing in the order $d\theta$ first and using $(1)$,
\[
  \int_{\mathbb{R}^n}\Bigl(\int_0^{2\pi} H(\theta,x)\,d\theta\Bigr) d\gamma_n(x)
    = \int_{\mathbb{R}^n} 8\,|v(x)| \, d\gamma_n(x) = 8M ,
\]
which is finite by Step 0.  Hence also
\begin{equation*}
  \int_0^{2\pi}\Bigl(\int_{\mathbb{R}^n} H(\theta,x)\, d\gamma_n(x)\Bigr) d\theta
    = 8M < \infty , \tag{2}
\end{equation*}
and Tonelli further gives that the inner integral
$\theta \mapsto \int H(\theta,x)\,d\gamma_n(x)$ is a measurable function of
$\theta$ on $[0,2\pi]$.

\emph{Step 6: the inner integral is $S_\theta$.}
Fix $\theta$.  By the description in the statement, $f_1(\theta)$ and
$f_2(\theta)$ are the $\gamma_n$-a.e.\ classes of the functions
$x \mapsto \operatorname{Re}(c_\theta v(x))$ and
$x \mapsto \operatorname{Im}(c_\theta v(x))$.  By
\noderef{ToL1} the $L^1(\gamma_n)$ class of a real $L^2(\gamma_n)$ function has
the same representative, so its $L^1$ norm is the integral of its absolute
value.  Therefore
\[
  \int_{\mathbb{R}^n} H(\theta,x)\,d\gamma_n(x)
    = \int\bigl|\operatorname{Re}(c_\theta v)\bigr| d\gamma_n
      + \int\bigl|\operatorname{Im}(c_\theta v)\bigr| d\gamma_n
    = \|f_1(\theta)\|_{L^1(\gamma_n)} + \|f_2(\theta)\|_{L^1(\gamma_n)}
    = S_\theta ,
\]
the first equality by linearity of the integral for the two nonnegative
measurable summands of $H$.

\emph{Step 7: conclusion.}
By Step 6 and Step 5, $\theta \mapsto S_\theta$ is measurable on $[0,2\pi]$ and
nonnegative with $\int_0^{2\pi} S_\theta\, d\theta = 8M < \infty$ by $(2)$; a
nonnegative measurable function with finite integral is integrable, which is
clause (1).  Dividing the identity $\int_0^{2\pi}S_\theta\,d\theta = 8M$ by
$2\pi > 0$,
\[
  \frac{1}{2\pi}\int_0^{2\pi} S_\theta\, d\theta = \frac{8M}{2\pi}
    = \frac{4}{\pi}M = \frac{4}{\pi}\int |v| \, d\gamma_n ,
\]
which is clause (2).
\end{proof}
